\honest{Newcomb's Problem and Regret of Rationality}{Eliezer Yudkowsky}{2008}

Omega, a superintelligence from another galaxy, correct on all 100 observed occasions, leaves two boxes and flies away. Box A holds \$1,000. Box B holds \$1,000,000 only if Omega predicted that you would take box B alone. The two-boxer says the boxes are already filled or empty, so taking both gains \$1,000 either way: \$1,000 instead of nothing, or \$1,001,000 instead of \$1,000,000. The one-boxer asks, ``If you're so rational, why ain'cha rich?'' The two-boxer answers that Omega chooses to reward irrational dispositions, too late to change. The literature is large, especially counting the Prisoner's Dilemma as a special case. \nb{Robert Nozick first published the problem in 1969. David Lewis took the taunt as the title of a 1981 paper for the two-boxers: ``The reason why we are not rich is that the riches were reserved for the irrational.''}

The dominant view, causal decision theory, says to take both boxes. \nb{The PhilPapers surveys agree: two boxes led one box 31.4\% to 21.3\% in 2009 and 39.0\% to 31.2\% in 2020, and decision theorists leaned further toward two boxes. Richard Jeffrey's evidential decision theory already recommended one box.} Even causal decision theorists agree that you should precommit to one box if you can, since that causes box B to be filled. In my field, self-modifying AI, an AI that two-boxes will modify itself to one-box if it foresees the problem; and one expecting Newcomblike problems of unknown form would make itself the kind of agent that generally does well on them.

I have a theory, but writing it up would take a small book, so I have shelved it. I know that verbal arguments for one-boxing are thought ``easy to come by,'' and a coherent theory hard. So I offer motivations instead. \nb{Yudkowsky later published timeless decision theory (2010) and, with Nate Soares, functional decision theory (2017).}

``Rational agents should WIN.'' Not selfishly or shortsightedly: if your utility function cares about others, win their happiness; about a million years hence, win the eon. Don't lose reasonably. Some think rationality predictably loses on some problems, the stereotype that Kirk beats Spock. A superbeing might reward a gene regardless of choice, or a particular algorithm, such as choosing the alphabetically last option, and not the same choice made another way. Omega rewards the choice to take one box whatever reasoning produced it, caring where we go, not how, so I deny that it rewards the irrational. \nb{The causal charge is that the choice itself is irrational. Lewis concluded that each side can consistently hold its view: ``So it's a standoff.''} Every rule is up for grabs except winning. For the same reason I will not bound my utility function to escape McGee's Dutch book over infinite times: for no finite N would I prefer an 80.0001\% chance of living N years to a 0.0001\% chance of a googolplex years and 80\% of living forever, so my utility is unbounded. Toss out the losing ritual; don't change the definition of winning, as if preferring \$1,000 to \$1,000,000 to make your ritual look good.

My causal decision theorist objects that one-boxing requires believing that my choice affects the box. My reply: I am a rationalist, so why care about being unreasonable? ``I can just\ldots\ take only box B.'' The point is to win, not to have an elegant theory of winning. Instead of asking whether ``reasonable'' agents get rich, look at the agents who leave with the money and work out from them what is reasonable; ``reasonable'' may just mean conforming to our current ritual. \nb{Lewis's charge is different: the evidential rule seeks good news about matters the choice cannot affect.}

James Joyce's causal decision theorist, Rachel, admits she wishes she were the refusing type, but says she is wishing ``for Irene's options, not sanctioning her choice.'' But she envies Irene ``only her choice,'' though you might rightly envy someone's genes; never envying another's mere choices is part of how I define rationality. My advice is: ``Just do the act you envy.'' Beware of defining the winner as anyone other than ``the agent who is currently smiling from on top of a giant heap of utility,'' or of calling cliff-refrainers unfairly advantaged over cliff-jumpers: ``Pay attention to the money!'' And if you would spend an extra hour trying to convince yourself that one box is rational, knowing you would then take it and find it full, that is an odd position, since rationality is meant to find the best choice, not a reason to believe in one. \nb{This restates the one-boxer's view. It does not answer Joyce's point that, once the prediction is made, Rachel's choice cannot change what is in the box.}

Perhaps two-boxing seems reasonable only while the money is not in front of you. If your daughter had a 90 per cent fatal disease, and box A held a serum with a 20 per cent chance of cure and box B might hold one with 95 per cent, or the boxes held asteroid deflectors working 10 and 100 per cent of the time, would you not wish that one-boxing were reasonable? ``Then maybe it's time to update your definition of reasonableness.'' Likewise, keeping separate track of the ``reasonable'' belief and the one likely true means one of them is wrong. \nb{Rachel already admits the wish, and denies that it shows the choice was mistaken.}

I admit this ``is not a knockdown criticism of causal decision theory.'' I use ``rational'' for my beliefs about accuracy and winning, not for verbal reasoning, certain success, proof, public demonstration or reasonableness; it currently means Bayescraft, and if Bayes ever lost systematically to an alternative by its mere decisions, it would go out the window. \nb{Both decision theories are Bayesian, so this does not choose between them.} Last, Musashi: ``The primary thing when you take a sword in your hands is your intention to cut the enemy, whatever the means''; if you think only of hitting, springing or parrying, you will not cut him.
